\section{Ring Switching}\label{sec:ring_switching}

\emph{Ring switching} is a reduction that turns a multilinear evaluation claim over a \emph{small}
coefficient ring $B$ into an evaluation claim over a \emph{large} extension $L$, paying only an
additive $O(1/\lvert L\rvert)$ soundness cost and \emph{without re-committing} over $L$. It is the
technique that lets a polynomial commitment scheme commit cheaply over a tiny ring (a binary-tower
field, $\mathbb{F}_2$, or a cyclotomic ring $R_q$) while running sum-check and the final opening over
a field large enough for soundness. It was introduced for binary towers in \cite{DP24}; the same
shape underlies Hachi's extension-field-to-cyclotomic-ring reduction \cite{NOZ26}.

ArkLib formalizes ring switching \emph{once}, as a generic oracle reduction parameterized by a
\emph{packing profile}, with the binary-tower (Binius) construction as one instance.

\subsection{The packing profile}

The data and algebraic laws that every ring-switching reduction abstracts over are collected into a
single structure.

\begin{definition}[Ring-switching profile]
  \lean{RingSwitching.RingSwitchingProfile}
  \label{def:ring_switching_profile}
  Fix commutative rings $B$ (small) and $L$ (large) with $L$ a $B$-algebra, and $\kappa \in
  \mathbb{N}$. A \emph{ring-switching profile} consists of:
  \begin{itemize}
    \item a rank-$2^\kappa$ $B$-basis $\beta = (\beta_u)_{u \in \{0,1\}^\kappa}$ of $L$;
    \item a commutative \emph{pack/trace carrier} $A$ that is an $L$-algebra (the type sent on the
      wire);
    \item two ring homomorphisms $\varphi_0, \varphi_1 : L \to A$ for evaluation-point data and
      polynomial coefficients, respectively;
    \item coordinate maps $\mathrm{decomposeColumns}, \mathrm{decomposeRows} : A \to (\{0,1\}^\kappa
      \to L)$;
    \item subject to the two \emph{reconstruction laws}
      \begin{equation}
        z = \sum_{u} \varphi_0(\beta_u)\cdot \varphi_1(\mathrm{decomposeRows}(z)_u).
        \label{eq:rs-rows}
      \end{equation}
      \begin{equation}
        z = \sum_{v} \varphi_0(\mathrm{decomposeColumns}(z)_v)\cdot \varphi_1(\beta_v).
        \label{eq:rs-cols}
      \end{equation}
      for all $z \in A$.
  \end{itemize}
\end{definition}

Laws \eqref{eq:rs-rows}--\eqref{eq:rs-cols} fix the coordinate convention of
\cite[\S2.5]{DP24}. They are data-layer reconstruction laws, not completeness or soundness
statements. Protocol correctness additionally needs identities connecting these coordinates to
packing, the honest folded element, and the batched multiplier. In a nontrivial carrier the laws
exclude an identically zero coordinate map; they do not by themselves establish a valid protocol
instance or non-vacuous security theorem.

\begin{definition}[Tensor-product profile]
  \lean{RingSwitching.tensorProductProfile}
  \uses{def:ring_switching_profile}
  \label{def:tensor_product_profile}
  For any fields $B = K$ and $L$ with a chosen degree-$2^\kappa$ extension basis, the profile takes the
  \emph{tensor carrier} $A := L \otimes_K L$, with $\varphi_0(\alpha) = \alpha \otimes 1$,
  $\varphi_1(\alpha) = 1 \otimes \alpha$, and the coordinate maps given by the $K$-basis coordinates
  for the left ($\beta.\mathrm{baseChange}\,L$) and right ($\mathrm{baseChangeRight}\,\beta$)
  $L$-module structures on $A$, giving columns and rows respectively. Thus
  $z = \sum_v z_v \otimes \beta_v = \sum_u \beta_u \otimes z_u$. Laws \eqref{eq:rs-rows}--\eqref{eq:rs-cols} hold by
  \texttt{Basis.sum\_repr} on these bases and are proven in ArkLib. Binius uses this constructor
  with its binary-tower basis.
\end{definition}

The Hachi construction \cite{NOZ26} is the intended second instance: $L = R_q$ a cyclotomic ring,
$A = R_q$ itself, $\varphi_0 = \mathrm{id}$, $\varphi_1 = \sigma_{-1}$ (the order-two automorphism),
and $\beta = \psi$ from its Theorem~2. Because $R_q$ is not an integral domain, the Schwartz--Zippel
soundness statement below does not apply to it; Hachi soundness is a separate argument.

\subsection{Packing the witness}

\begin{definition}[Multilinear packing]
  \lean{RingSwitching.packMLE}
  \uses{def:ring_switching_profile}
  \label{def:pack_mle}
  With $\ell = \ell' + \kappa$, a small-field multilinear $t$ in $\ell$ variables is packed into a
  large-field multilinear $t'$ in $\ell'$ variables by bundling each block of $2^\kappa$ coefficients
  into one $L$-element:
  \[
    t'(w) \;=\; \sum_{v \in \{0,1\}^\kappa} t(v, w)\,\beta_v, \qquad w \in \{0,1\}^{\ell'}.
  \]
\end{definition}

\subsection{The reduction}

The full reduction is a sequential composition of three phases, ending in a call to an underlying
large-field multilinear IOPCS over $L$.

\begin{enumerate}
  \item \textbf{Batching} (\lean{RingSwitching.BatchingPhase.batchingOracleReduction}). The prover
    sends the folded element $\hat s \in A$. The verifier checks the original claim against its
    column-decomposition, $s = \sum_v \tilde{\mathrm{eq}}(v, r_{0..\kappa-1})\cdot
    \mathrm{decomposeColumns}(\hat s)_v$, then sends a batching challenge $r'' \in L^\kappa$,
    collapsing the $\kappa$ packed coordinates into a single sum-check target $s_0 = \sum_u
    \tilde{\mathrm{eq}}(u, r'')\cdot \mathrm{decomposeRows}(\hat s)_u$. (Soundness loss
    $\kappa/\lvert L\rvert$.)
  \item \textbf{Sum-check} (\lean{RingSwitching.SumcheckPhase.coreInteractionOracleReduction}). A
    degree-$2$ multilinear sum-check over $L$ on $H = A_{\mathrm{MLE}} \cdot t'$, over $\ell'$ rounds
    plus a final round, reducing the sum claim to a single evaluation of $t'$ at a random point.
    (Loss $2/\lvert L\rvert$ per round, $1/\lvert L\rvert$ final.)
  \item \textbf{Opening}. The residual evaluation claim about $t'$ over $L$ is discharged by the
    underlying large-field multilinear IOPCS.
\end{enumerate}

\begin{definition}[Full ring-switching reduction]
  \lean{RingSwitching.FullRingSwitching.fullOracleReduction}
  \uses{def:ring_switching_profile, def:pack_mle}
  \label{def:full_ring_switching}
  The composition of the batching phase, the core sum-check interaction, and the large-field IOPCS
  invocation.
\end{definition}

\subsection{Security}

\begin{theorem}[Perfect completeness target]
  \lean{RingSwitching.FullRingSwitching.fullOracleReduction_perfectCompleteness}
  \uses{def:full_ring_switching}
  \label{thm:ring_switching_completeness}
  The full ring-switching oracle reduction is perfectly complete: an honest prover holding a witness
  in the input relation always produces a transcript accepted into the output relation. The honest
  prover's folded $\hat s$ must satisfy the verifier's column check through the packing and
  folded-evaluation identities, in addition to reconstruction law \eqref{eq:rs-cols}. These
  instance-specific identities remain proof obligations in the generic development.
\end{theorem}

\begin{theorem}[Round-by-round knowledge soundness target]
  \lean{RingSwitching.FullRingSwitching.fullOracleVerifier_rbrKnowledgeSoundness}
  \uses{def:full_ring_switching}
  \label{thm:ring_switching_soundness}
  When $L$ is an integral domain, the full reduction is round-by-round knowledge sound, with the
  additive error
  \[
    \frac{\kappa}{\lvert L\rvert} \;+\; \sum_{\text{rounds}} \frac{2}{\lvert L\rvert} \;+\;
    \frac{1}{\lvert L\rvert} \;+\; \varepsilon_{\mathrm{IOPCS}},
  \]
  matching \cite[\S3.1--3.2]{DP24}. The integral-domain hypothesis is exactly the Schwartz--Zippel
  ingredient (a nonzero degree-$d$ univariate over $L$ has at most $d$ roots).
\end{theorem}

\noindent The leaf completeness and soundness proofs are in progress. The displayed security
statements describe targets, not established end-to-end Binius guarantees. A correct instance
requires the packing, folded-evaluation, batching, and final-multiplier identities as well as the
data-layer reconstruction laws. The concrete tensor instance and its coordinate regressions do
not discharge those full protocol obligations.
